% Generated draft: requires mathematical and source review.
\documentclass[11pt]{article}
\usepackage[margin=25mm]{geometry}
\usepackage{fontspec}
\setmainfont{Latin Modern Roman}
\usepackage{amsmath,amssymb,mathtools,unicode-math}
\setmathfont{Latin Modern Math}
\usepackage{enumitem,needspace,xurl,hyperref}
\setlength{\parindent}{0pt}
\setlength{\parskip}{5pt}
\setlength{\emergencystretch}{3em}
\newcommand{\N}{\mathbb{N}}
\newcommand{\Z}{\mathbb{Z}}
\newcommand{\Q}{\mathbb{Q}}
\newcommand{\R}{\mathbb{R}}
\newcommand{\C}{\mathbb{C}}
\newcommand{\F}{\mathbb{F}}
\newcommand{\A}{\mathbb{A}}
\newcommand{\PP}{\mathbb{P}}
\newcommand{\E}{\mathbb{E}}
\newcommand{\eps}{\varepsilon}
\newcommand{\dd}{\,\mathrm{d}}
\begin{document}
\section*{UnsolvedMath ID: 701367}
\textbf{Source problem code:} catalogue entry
\subsection{Definitions and mathematical statement}
The supplied record depends on Problem 17.64, whose statement is not supplied. In particular, the meanings of \(p\), \(J\), and a \(g(p)\)-approximation to \(J\), as well as the admissible functions \(g\), are not specified. A complete mathematical formulation therefore cannot be recovered from this record alone.

For the sole purpose of displaying the requested cardinality conditions, write formally
\[
\mathcal I_{g,p}(J)=\{[X]_{\cong}:X\text{ is a }g(p)\text{-approximation to }J\},
\]
where \([X]_{\cong}\) denotes an isomorphism class. This notation does not supply the missing definition of approximation or specify the objects and isomorphisms involved.

Assuming that Problem 17.64 has an affirmative answer, does it follow that there is some function \(g(p)\) for which the number of isomorphism classes of \(g(p)\)-approximations to \(J\) is:
\begin{enumerate}
\item[(a)] countable;
\item[(b)] exactly one, that is, \(\lvert\mathcal I_{g,p}(J)\rvert=1\)?
\end{enumerate}
These are separate questions under the same stated premise. The domains, admissibility conditions, and quantification over \(p\) and \(J\) must be recovered from the missing context; no uniformity in either parameter is imposed here. The source does not explicitly distinguish ``at most countable'' from ``countably infinite'' in part (a).

The supplied source additionally states: if, for some \(g(p)\), there are only finitely many isomorphism classes, then for some \(h(p)\) there is only one. This is retained as a source qualification, with the same unresolved parameter domains and approximation definitions; it is not independently established here.
\subsection{Short English statement}
C. R. Leedham-Green asks whether an affirmative answer to Problem 17.64 implies that a suitable approximation parameter \(g(p)\) yields (a) countably many or (b) exactly one isomorphism class of approximations to \(J\). The source also states that finiteness for one such function implies uniqueness for another. Problem 17.64 and the definitions needed to make these questions precise are absent from the supplied record.
\subsection{Sources}
\begin{itemize}
\item Local UnsolvedMath registry, permanent ID \(701367\), exact source code \texttt{catalogue entry}, title \emph{Kourovka Notebook Problem 17.65}; supplied upstream revision \texttt{d371ee917ac9647fbb06cd6fc529deec6f3dda82}.
\item Supplied bibliographic attribution: \emph{The Kourovka Notebook (Unsolved Problems in Group Theory)}, 17th issue (2010), Problem 17.65; also identified with the 21st issue edition (2026), \texttt{arXiv:1401.0300v48}. Proposed by C. R. Leedham-Green in 2010. The supplied collection description names E. I. Khukhro and V. D. Mazurov as editors.
\item The registry labels this problem \emph{open}. This is a supplied classification, not an independently verified assessment. The cited notebook text and the prerequisite Problem 17.64 have not been independently consulted for this draft.
\end{itemize}
\end{document}
